\chapter{Weak approximation, divisors, and genus}
\label{par-chap-divisors}

\section{Weak approximation}
\begin{theorem}[Weak Approximation Theorem]
\label{par-thm-weak-approximation}
\dcref{1.9}
\uses{par-def-place,par-def-place-valuation,par-def-function-field}
\source{papaioannou-algebraic-rr:page-0006}
Given pairwise distinct places $P_1,\ldots,P_n$, elements $x_1,\ldots,x_n\in\mathbb F$,
and integers $r_1,\ldots,r_n$, there is $x\in\mathbb F$ with
\[
v_{P_i}(x-x_i)=r_i\qquad(1\leq i\leq n).
\]
\end{theorem}
\begin{proof}
\notready
\end{proof}

\begin{remark}[Valuation absolute values]
\label{par-rem-weak-approximation-absolute-value}
\uses{par-thm-weak-approximation,par-def-valuation}
\source{papaioannou-algebraic-rr:page-0006}
The valuation $v_P$ induces $|z|_P=2^{-v_P(z)}$ for $z\neq0$ and $|0|_P=0$.
Weak approximation says that finitely many prescribed elements can be
simultaneously approximated arbitrarily well by one element.
\end{remark}

\begin{corollary}[Infinitely many places]
\label{par-cor-infinitely-many-places}
\dcref{1.10}
\uses{par-thm-weak-approximation,par-cor-zero-pole}
\source{papaioannou-algebraic-rr:page-0006}
Every function field has infinitely many places.
\end{corollary}
\begin{proof}
If the places were $P_1,\ldots,P_n$, weak approximation would give a nonzero
element with positive valuation at each $P_i$. It would be transcendental but
have no pole, contradicting Corollary~1.7.
\end{proof}

\begin{theorem}[Weighted zero bound]
\label{par-thm-weighted-zero-bound}
\dcref{1.11}
\uses{par-def-zero-pole,par-thm-place-degree}
\source{papaioannou-algebraic-rr:page-0006}
If $P_1,\ldots,P_r$ are zeros of $x\in\mathbb F$, then
\[
\sum_{i=1}^r v_{P_i}(x)\deg P_i\leq[\mathbb F:\mathbb K(x_i)].
\]
The field in the final bound is $\mathbb K(x)$.
\end{theorem}
\begin{proof}
\notready
\end{proof}

\begin{corollary}[Finite support of zeros and poles]
\label{par-cor-finite-zeros-poles}
\dcref{1.12}
\uses{par-thm-weighted-zero-bound,par-def-zero-pole}
\source{papaioannou-algebraic-rr:page-0006}
Every $x\in\mathbb F^*$ has only finitely many zeros and poles.
\end{corollary}
\begin{proof}
Constants have neither. For a nonconstant element, the preceding bound applies
to its zeros, and applying it to $x^{-1}$ gives finiteness of its poles.
\end{proof}

\section{Divisors and the genus}
The constant field is henceforth assumed algebraically closed in $\mathbb F$.

\begin{definition}[Divisor and support]
\label{par-def-divisor}
\dcref{1.6}
\uses{par-def-place,par-cor-finite-zeros-poles}
\source{papaioannou-algebraic-rr:page-0006}
$\Div(\mathbb F)$ is the free abelian group generated by
$\mathbf P_{\mathbb F}$. A divisor is a finite formal sum
\[
D=\sum_{P\in\mathbf P_{\mathbb F}}n_PP,\qquad n_P\in\mathbb Z,
\]
and $\Supp(D)=\{P:n_P\neq0\}$.
\end{definition}

\begin{definition}[Order, effectivity, and degree]
\label{par-def-divisor-degree}
\uses{par-def-divisor,par-def-place-valuation}
\source{papaioannou-algebraic-rr:page-0007}
Write $v_P(D)=n_P$. Define $D_1\leq D_2$ when
$v_P(D_1)\leq v_P(D_2)$ for all $P$; $D$ is effective when $D\geq0$. The degree is
\[
\deg D=\sum_{P\in\mathbf P_{\mathbb F}}v_P(D)\deg P,
\]
a homomorphism $\Div(\mathbb F)\to\mathbb Z$.
\end{definition}

\begin{definition}[Zero, pole, and principal divisors]
\label{par-def-principal-divisor}
\dcref{1.7}
\uses{par-def-divisor,par-def-zero-pole}
\source{papaioannou-algebraic-rr:page-0007}
For $x\in\mathbb F^*$, with zero set $Z$ and pole set $N$, set
\[
(x)_0=\sum_{P\in Z}v_P(x)P,\qquad
(x)_\infty=\sum_{P\in N}(-v_P(x))P,\qquad
(x)=(x)_0-(x)_\infty=\sum_Pv_P(x)P.
\]
\end{definition}

\begin{definition}[Divisor class group]
\label{par-def-class-group}
\dcref{1.8}
\uses{par-def-divisor,par-def-principal-divisor}
\source{papaioannou-algebraic-rr:page-0007}
The subgroup $\Princ(\mathbb F)=\{(x):x\in\mathbb F^*\}$ is the group of
principal divisors, with $(xy)=(x)+(y)$. The class group is
\[
\Class(\mathbb F)=\Div(\mathbb F)/\Princ(\mathbb F).
\]
Equality of classes defines $D\sim D'$, equivalently $D-D'=(x)$ for some $x$.
\end{definition}

\begin{definition}[Riemann--Roch space]
\label{par-def-L-space}
\dcref{1.9}
\uses{par-def-divisor,par-def-principal-divisor}
\source{papaioannou-algebraic-rr:page-0007}
For a divisor $D$,
\[
L(D)=\{x\in\mathbb F^*:(x)\geq-D\}\cup\{0\}.
\]
This is a $\mathbb K$-vector space. Explicitly, if $v_P(D)>0$ precisely on a finite set
$I$ and $v_Q(D)<0$ precisely on a finite set $J$, then $L(D)$ consists of
functions with zeroes of order at least $-v_Q(D)$ at $Q\in J$ and poles of order
at most $v_P(D)$ at $P\in I$. If $D\sim D'$ via $D=D'+(z)$, multiplication by
$z$ identifies $L(D)$ and $L(D')$.
\end{definition}

\begin{definition}[Genus]
\label{par-def-genus}
\dcref{1.10}
\uses{par-def-divisor,par-def-L-space}
\source{papaioannou-algebraic-rr:page-0007}
The genus of $\mathbb F/\mathbb K$ is
\[
g=\sup\{\deg D-\dim_{\mathbb K}L(D)+1:D\in\Div(\mathbb F)\}.
\]
\end{definition}

\begin{remark}[Supremum convention]
\label{par-rem-genus-supremum}
\uses{par-def-genus,par-def-L-space}
\source{papaioannou-algebraic-rr:page-0007}
The quantity $\deg D-\dim D$ is bounded above, so the supremum is finite.
\end{remark}

\begin{lemma}[Riemann inequality and eventual equality]
\label{par-lem-riemann-inequality}
\uses{par-def-genus,par-def-L-space,par-def-divisor-degree}
\source{papaioannou-algebraic-rr:page-0008}
For every $D$, $g\geq\deg D-\dim D+1$. There is an integer $c$ such that
$\dim D=\deg D+1-g$ whenever $\deg D\geq c$.
\end{lemma}
\begin{proof}
Choose $D_0$ attaining the supremum and put $c=\deg D_0+g$. For
$\deg D\geq c$, Riemann's inequality gives
$\dim(D-D_0)\geq\deg(D-D_0)+1-g\geq1$. Choose
$0\neq z\in L(D-D_0)$ and set $D'=D+(z)\geq D_0$. Then
$\deg D-\dim D=\deg D'-\dim D'\geq g-1$, giving the reverse inequality to
Riemann's inequality and hence equality.
\end{proof}

\begin{lemma}[Inclusion and dimension bound]
\label{par-lem-L-inclusion-bound}
\dcref{1.13}
\uses{par-def-L-space,par-def-divisor-degree,par-def-place-valuation}
\source{papaioannou-algebraic-rr:page-0008}
If $D\leq E$, then $L(D)\subseteq L(E)$ and
\[
\dim L(E)/L(D)\leq\deg E-\deg D.
\]
\end{lemma}
\begin{proof}
Reduce inductively to $E=D+P$. Choose $t$ with
$v_P(t)=v_P(D)+1$. The map
$\phi:L(E)\to\mathbb F_P$, $x\mapsto(xt)(P)$, has kernel $L(D)$, hence induces
an injection $L(E)/L(D)\hookrightarrow\mathbb F_P$. Its dimension is at most
$\deg P=\deg E-\deg D$.
\end{proof}

\begin{lemma}[Finite dimensionality of $L(D)$]
\label{par-lem-L-finite}
\dcref{1.14}
\uses{par-lem-L-inclusion-bound,par-def-L-space,par-def-divisor-degree}
\source{papaioannou-algebraic-rr:page-0008}
Writing $D=D_+-D_-$ with $D_+,D_-\geq0$, one has
$\dim L(D)\leq\deg D_++1$; in particular $L(D)$ is finite-dimensional.
\end{lemma}
\begin{proof}
$L(D)\subseteq L(D_+)$, and the preceding lemma gives
$\dim(L(D_+)/L(0))\leq\deg D_+$. Since $L(0)=\mathbb K$ under the full-constant-
field assumption, $\dim L(D_+)\leq\deg D_++1$.
\end{proof}

\begin{theorem}[Principal divisors have degree zero]
\label{par-thm-principal-degree-zero}
\dcref{1.15}
\uses{par-thm-weighted-zero-bound,par-lem-L-finite,par-def-principal-divisor}
\source{papaioannou-algebraic-rr:page-0008}
For every $x\in\mathbb F^*$,
\[
\deg(x)_0=\deg(x)_\infty=[\mathbb F:\mathbb K(x)].
\]
\end{theorem}
\begin{proof}
Let $n=[\mathbb F:\mathbb K(x)]$ and $D=(x)_\infty$. The weighted zero bound
applied to $x^{-1}$ gives $\deg D\leq n$. Choose a basis $u_1,\ldots,u_n$ of
$\mathbb F/\mathbb K(x)$ and $C\geq0$ with $(u_i)\geq-C$. Then
$x^ju_i\in L(kD+C)$ for $0\leq j\leq k$, so
$n(k+1)\leq\dim(kD+C)\leq k\deg D+\deg C+1$. Letting $k$ grow gives
$\deg D\geq n$. Thus $\deg(x)_\infty=n$; apply the same argument to $x^{-1}$.
\end{proof}

\begin{corollary}[Equivalent divisors and negative degree]
\label{par-cor-equivalent-divisors}
\dcref{1.16}
\uses{par-thm-principal-degree-zero,par-def-class-group,par-def-L-space}
\source{papaioannou-algebraic-rr:page-0009}
\begin{enumerate}
\item Equivalent divisors have equal degree and dimension.
\item If $\deg D<0$, then $\dim D=0$.
\item For $\deg D=0$, the following are equivalent: (i) $D$ is principal;
(ii) $\dim D\geq1$; (iii) $\dim D=1$.
\end{enumerate}
\end{corollary}
\begin{proof}
\notready
\end{proof}

\begin{theorem}[Finite genus]
\label{par-thm-finite-genus}
\dcref{1.17}
\uses{par-lem-L-inclusion-bound,par-lem-L-finite,par-def-genus,par-thm-principal-degree-zero}
\source{papaioannou-algebraic-rr:page-0009}
There is $\gamma\in\mathbb Z$, depending only on $\mathbb F/\mathbb K$, such that
\[
\deg D-\dim D\leq\gamma\qquad(D\in\Div(\mathbb F)).
\]
\end{theorem}
\begin{proof}
Fix nonconstant $x$ and $B=(x)_\infty$. With $C\geq0$ as in the proof of
Theorem~1.15, $\dim(kB)\geq\deg(kB)+[\mathbb F:\mathbb K(x)]-\deg C$.
Thus the asserted bound holds on multiples of $B$. For arbitrary $A$, choose
effective $A_1\geq A$ and large $k$ with
$\dim(kB-A_1)>0$. A nonzero $z\in L(kB-A_1)$ gives
$D=A_1-(z)_0\sim A_1$ and $D\leq kB$, so monotonicity from Lemma~1.13 transfers
the same bound to $A$.
\end{proof}
